\begin{proposition}
    \label{prop:periodicidad}
    Sea $\alpha \in \mathbb{R}$.
    \begin{enumerate}
        \item Si $\alpha \in \mathbb{Q}$, entonces la sucesión $(\partfrac{n\alpha})_{n\ge1}$ es periódica
        y tiene un número finito de términos distintos.
        \item Si $\alpha \notin \mathbb{Q}$, entonces la sucesión $(\partfrac{n\alpha})_{n\ge1}$ tiene
        infinitos términos distintos.
    \end{enumerate}
\end{proposition}

\begin{proof}
    \textit{1. \quad} Sea $\alpha = p/q$ con $p, q \in \mathbb{Z}$ y $q \neq 0$. Entonces,
    \[
        \partfrac{p/q}, \,
        \partfrac{2p/q}, \,
        \ldots, \,
        \partfrac{(q-1)p/q}, \,
        \partfrac{qp/q} = \partfrac{p} = 0 
    \]
    A partir de aquí, la sucesión se repite,
    \[
        \partfrac{(q+1)p/q} = \partfrac{1 + p/q} = \partfrac{p/q}.
    \]
    \textit{2. \quad} Sea $\alpha \notin \mathbb{Q}$ y supongamos que $\partfrac{n\alpha} = \partfrac{m\alpha}$ 
    para ciertos $n, m \in \mathbb{N}$ con $n \neq m$. Tenemos que
    \begin{align*}
        \partfrac{n\alpha} = \partfrac{m\alpha}
        &\iff n\alpha - \floor{n\alpha} = m\alpha - \floor{m\alpha} \\
        &\iff (n-m)\alpha = \floor{n\alpha} - \floor{m\alpha} \in \mathbb{Z}
    \end{align*}
    lo que implicaría que $\alpha \in \mathbb{Q}$. Contradicción \#.
\end{proof}
